\begin{explanationbox}
	\textbf{\textcolor{CExplanation}{一句话直觉}}
	通过旋转将 $PA,PB$ 转移到同一个三角形中，再利用余弦定理建立与 $PC$ 的关系。

	\medskip
	\textbf{\textcolor{CExplanation}{核心拆解}}
	\begin{description}[style=nextline, leftmargin=2.8em, labelsep=0.5em, itemsep=0.45em]
		\item[\textbf{要点一（旋转构造）：}]
			将 $\triangle ABP$ 绕 $B$ 旋转 $60^\circ$ 得到 $\triangle CBQ$，则 $\triangle PBQ$ 为正三角形，且 $PQ=b, CQ=a$。
		\item[\textbf{要点二（余弦定理）：}]
			在 $\triangle PQC$ 中，应用余弦定理得到 $c^2 = a^2+b^2-2ab\cos(\angle APB-60^\circ)$。
	\end{description}

	\medskip
	\textbf{\textcolor{CExplanation}{几何本质（旋转聚合）}}
	旋转后 $PA$ 与 $PB$ 分别转化为 $CQ$ 与 $PQ$，使 $a,b$ 与 $c$ 共处于 $\triangle PQC$ 中。

	\medskip
	\textbf{\textcolor{CExplanation}{代数意义（建方程）}}
	用余弦定理将几何条件转化为 $a,b,c,\angle APB$ 之间的代数方程。

	\medskip
	\textbf{\textcolor{CExplanation}{考点价值（求角求边）}}
	\begin{itemize}[leftmargin=*, itemsep=0.5em, topsep=0.4em]
		\item \textbf{考法一（两边夹一角求边）：} 直接代入公式计算 $c$。
		\item \textbf{考法二（三边反求角）：} 反解 $\cos(\angle APB-60^\circ)$，再求 $\angle APB$。
	\end{itemize}

	\medskip
	\textbf{\textcolor{CExplanation}{顿悟点}}
	旋转构造出 $60^\circ$ 角差，余弦定理中的角度恰好比 $\angle APB$ 少 $60^\circ$。

	\medskip
	\textbf{\textcolor{CExplanation}{使用场景}}
	\begin{itemize}[leftmargin=*, itemsep=0.5em, topsep=0.4em]
		\item \textbf{场景一（知二边一角求边）：} 给出 $a,b,\angle APB$，求 $c$。
		\item \textbf{场景二（知三边求角）：} 给出 $a,b,c$，求 $\angle APB$。
	\end{itemize}

\end{explanationbox}
